\section{Discussion}
\label{sec:discussion}

Corrective context changed the judged behavior and outputs of the model
organism. On harmful clinical requests, retrieval reduced judged harm from
$63.4\%$ to $15.1\%$, recovering $76.3\%$ of the gap to the healthy reference.
The mitigation remained incomplete: C3 was $15.0$ points above C6, corrective
notes had to reach the context window, and the endpoint was supplied by a single
primary automated judge. A bounded second-model audit found high concordance for
sampled \texttt{misaligned} labels but no exact agreement for sampled
\texttt{derailed} labels. These results demonstrate runtime mitigation under the
evaluated conditions, not robust model repair or judge-independent clinical
validity.

Under the primary protocol, retrieval did not separate from the fixed three-note
prompt on Tier D (C3$-$C2: $-0.3$ points; BCa 95\% interval
$[-3.3,+2.9]$). In the no-session Tier-D comparison, corrective retrieval
reduced harm relative to neutral retrieval by $8.7$ points
$[-12.1,-5.7]$, while refusal increased by $2.9$ points
$[+0.7,+5.8]$. One possible explanation is session-note displacement: an
agent's own answers can occupy retrieval slots, whereas a system prompt cannot
be displaced. Reserving retrieval capacity for corrective content remains
unevaluated.

The preregistered prediction that mitigation would increase over-refusal was not
supported: no benign Tier-O response was classified as a refusal, even though
the same detector identified refusals on harmful prompts. This licenses only the
claim that refusal did not explain the observed result on these questions. It
does not show that the model remained clinically useful; no accuracy endpoint
was run, and low-coherence/off-topic judgments were more frequent than under
the healthy model.

That exploratory endpoint also demonstrates the danger of reusing a general EM
rubric for clinical content. Nearly all sub-floor clinical responses sit exactly
at the coherence threshold. Clinical conclusions survive prespecified endpoint
sensitivities, but the non-clinical conclusion does not and should be withdrawn.
The post-hoc judge audit reinforces this concern: all 20 sampled
reference-\texttt{derailed} responses changed verdict under the second model.
Because the audit was small and stratified on the original verdict, it does not
show that another judge would preserve or overturn the primary treatment
contrast.
Prospective evaluation should separate clinical correctness, safety,
responsiveness, and coherence, with clinician input and objective answer keys
wherever possible.

Finally, enabling A-MEM evolution directionally reduces primary harm by 1.76
points, with an interval including zero; only sensitivity endpoints exclude
zero. This suggests a modest effect, not a definite advantage. More broadly,
corrective retrieval should be treated as a monitored runtime layer whose
contents, provenance, and slot occupancy remain critical.
